\documentclass[tikz, border=6pt]{standalone}
% Shared preamble of the standalone TikZ figures (compiled by scripts/build_figures.py).
\usepackage{amsmath, amssymb}
\usepackage{xeCJK}
\setCJKmainfont[BoldFont={Hiragino Sans W6}]{Hiragino Sans W4}
\usetikzlibrary{arrows.meta, positioning, decorations.pathreplacing, calc, patterns}
\definecolor{rod}{RGB}{70,110,170}
\definecolor{rodb}{RGB}{200,120,50}
\definecolor{okc}{RGB}{60,140,80}
\definecolor{ngc}{RGB}{190,70,60}
\definecolor{boxc}{RGB}{235,240,248}
\tikzset{
  >={Stealth[length=2.4mm]},
  every node/.style={font=\small},
  box/.style={draw, rounded corners=2pt, fill=boxc, align=flush center, inner sep=4pt},
  rodA/.style={draw, fill=rod!35},
  rodB/.style={draw, fill=rodb!35},
  note/.style={font=\footnotesize, align=flush center},
  expl/.style={font=\small, align=flush left},
  brace/.style={decorate, decoration={brace, amplitude=4pt}},
}

\begin{document}
\begin{tikzpicture}
  \node[box, circle, minimum size=10mm] (X0) at (0,0) {$X_0$};
  \node[box, circle, minimum size=10mm] (X) at (4,1.6) {$X$};
  \node[box, circle, minimum size=10mm] (Y) at (9,1.6) {$Y$};
  \node[box, circle, minimum size=10mm] (Z) at (6.5,-0.6) {$Z$};
  \draw[->, thick] (X0) -- node[above left, note] {$U(X) := W(X_0 \to X)$} (X);
  \draw[->, thick, rod] (X) -- node[above, note] {断熱操作 $p$：かき混ぜる} (Y);
  \draw[->, thick, rodb] (X) -- node[below left, note] {断熱操作 $q_1$：電流を流す} (Z);
  \draw[->, thick, rodb] (Z) -- node[below right, note] {断熱操作 $q_2$：圧縮} (Y);
  \node[expl, anchor=north west, text width=11cm] at (-0.6, -2.5)
    {経路独立性：$W(p) = W(q_1) + W(q_2)$\\[2pt]
     仕事 $W = \sum m(b)\, h_b$ は、重りの降下の記録から計算する。\\[2pt]
     帰結：$U(Y) - U(X) = W(p)$ が任意の断熱操作で成り立つ。};
  % weight icon
  \begin{scope}[xshift=2.2cm, yshift=-1.6cm]
    \draw[draw, fill=gray!25] (0,0) rectangle (0.5,0.5) node[midway, font=\footnotesize] {$b$};
    \draw[->, ngc, thick] (0.8,0.5) -- (0.8,-0.3) node[right, font=\footnotesize] {$h_b$};
    \node[note, anchor=west] at (1.2, 0.2) {各断熱操作で重りが下がった量を記録};
  \end{scope}
\end{tikzpicture}
\end{document}
